% !TeX program = pdflatex
\documentclass[aspectratio=169,xcolor={dvipsnames,table}]{beamer}

% Preámbulo modular para presentaciones Beamer (Modelación Estadística)
% Diseño y paleta institucional del Tecnológico de Monterrey

\documentclass[aspectratio=169,xcolor={dvipsnames,table}]{beamer}

\usepackage[utf8]{inputenc}
\usepackage[spanish,mexico]{babel}
\usepackage{amsmath,amssymb,amsfonts,mathtools}
\usepackage{graphicx}
\usepackage{listings}
\usepackage{booktabs}
\usepackage{tikz}
\usetikzlibrary{matrix,arrows,decorations.pathmorphing,shapes.geometric,positioning}

% Paleta Institucional Tecnológico de Monterrey
\definecolor{TecAzul}{HTML}{0039A6}       % Azul TEC: color primario institucional
\definecolor{TecAzulOscuro}{HTML}{1A2E51} % Azul marino oscuro
\definecolor{TecRojo}{HTML}{EC2661}       % Rojo de acento (paleta miTec)
\definecolor{TecGris}{HTML}{646464}       % Gris neutro para comentarios y subtítulos
\definecolor{TecFondoBloque}{HTML}{F4F6F9}% Fondo suave para bloques

% Configuración de Tema Beamer
\usetheme{Boadilla}
\usecolortheme[named=TecAzulOscuro]{structure}
\setbeamercolor{alerted text}{fg=TecRojo}
\setbeamercolor{palette primary}{bg=TecAzulOscuro,fg=white}
\setbeamercolor{palette secondary}{bg=TecAzul,fg=white}
\setbeamercolor{palette tertiary}{bg=TecAzulOscuro!80!black,fg=white}
\setbeamercolor{title}{fg=TecAzulOscuro,bg=TecFondoBloque}
\setbeamercolor{frametitle}{fg=TecAzulOscuro,bg=TecFondoBloque}

% Bloques personalizados elegantes
\setbeamercolor{block title}{bg=TecAzulOscuro,fg=white}
\setbeamercolor{block body}{bg=TecFondoBloque,fg=black}
\setbeamercolor{block title alerted}{bg=TecRojo,fg=white}
\setbeamercolor{block body alerted}{bg=TecRojo!8,fg=black}
\setbeamercolor{block title example}{bg=TecAzul,fg=white}
\setbeamercolor{block body example}{bg=TecAzul!8,fg=black}

% Redondear bordes y añadir sombra sutil
\setbeamertemplate{blocks}[rounded][shadow=true]
\setbeamertemplate{navigation symbols}{}

% Pie de página limpio con número de diapositiva
\setbeamertemplate{footline}{
  \leavevmode%
  \hbox{%
  \begin{beamercolorbox}[wd=.7\paperwidth,ht=2.25ex,dp=1ex,left]{author in head/foot}%
    \hspace*{2ex}\insertshortauthor~~(\insertshortinstitute) \hfill \insertshorttitle
  \end{beamercolorbox}%
  \begin{beamercolorbox}[wd=.3\paperwidth,ht=2.25ex,dp=1ex,right]{date in head/foot}%
    \insertshortdate{}\hspace*{2em}
    \insertframenumber{} / \inserttotalframenumber\hspace*{2ex}
  \end{beamercolorbox}}%
  \vskip0pt%
}

% Configuración de Listings de Python para visualización pedagógica de código
\lstset{
    language=Python,
    basicstyle=\ttfamily\footnotesize,
    keywordstyle=\color{TecAzulOscuro}\bfseries,
    stringstyle=\color{TecRojo},
    commentstyle=\color{TecGris}\itshape,
    showstringspaces=false,
    breaklines=true,
    frame=lines,
    rulecolor=\color{TecAzulOscuro!30},
    backgroundcolor=\color{TecFondoBloque},
    aboveskip=0.5em,
    belowskip=0.5em,
    literate={á}{{\'a}}1 {é}{{\'e}}1 {í}{{\'i}}1 {ó}{{\'o}}1 {ú}{{\'u}}1
             {Á}{{\'A}}1 {É}{{\'E}}1 {Í}{{\'I}}1 {Ó}{{\'O}}1 {Ú}{{\'U}}1
             {ñ}{{\~n}}1 {Ñ}{{\~N}}1 {¿}{{?`}}1 {¡}{{!`}}1
}

% Metadatos institucionales por defecto
\title[Modelación Estadística]{Modelación Estadística para Ciencia de Datos e Ingeniería}
\author[J. Castillo Colmenares]{Juliho David Castillo Colmenares}
\institute[Tec de Monterrey]{Tecnológico de Monterrey \\ Escuela de Ingeniería y Ciencias}
\date{\today}

% Comandos y macros estadísticas compartidas para las presentaciones Beamer (Metropolis)

% Conjuntos numéricos básicos
\newcommand{\R}{\mathbb{R}}
\newcommand{\N}{\mathbb{N}}
\newcommand{\Q}{\mathbb{Q}}
\newcommand{\Z}{\mathbb{Z}}
\newcommand{\set}[1]{\left\{ #1 \right\}}
\newcommand{\abs}[1]{\left|#1\right|}
\newcommand{\norm}[1]{\left\| #1 \right\|}

% Comandos estadísticos y probabilísticos del curso
\newcommand{\Var}{\operatorname{Var}}
\newcommand{\cov}{\operatorname{Cov}}
\newcommand{\corr}{\rho}
\newcommand{\comb}[2]{\begin{pmatrix} #1 \\ #2 \end{pmatrix}}
\newcommand{\s}{\sigma}
\newcommand{\particion}{\mathfrak{P}}
\newcommand{\card}[1]{\#\left( #1 \right)}
\newcommand{\seq}[3]{\set{#1}_{#2}^{#3}}
\newcommand{\E}{\mathbb{E}}
\newcommand{\Prob}{\mathbb{P}}

% Entornos pedagógicos y cajas en español académico natural
\newenvironment{discusion}{%
  \begin{alertblock}{Pregunta de Discusión}%
}{%
  \end{alertblock}%
}

\newenvironment{reflexion}{%
  \begin{alertblock}{Reflexión para el Aula}%
}{%
  \end{alertblock}%
}

\newenvironment{paradoja}{%
  \begin{alertblock}{Paradoja de Exploración}%
}{%
  \end{alertblock}%
}

\newenvironment{motivacion}{%
  \begin{exampleblock}{Motivación: Ciencia de Datos en la Práctica}%
}{%
  \end{exampleblock}%
}

\newenvironment{retoclase}{%
  \begin{block}{Reto Rápido en Equipo}%
}{%
  \end{block}%
}


\title{Distribución Hipergeométrica}
\subtitle{Sección 03.06 --- Muestreo Sin Reemplazo y FPCF}
\author[J. Castillo Colmenares]{Juliho Castillo Colmenares}
\institute[Tec de Monterrey]{Tecnológico de Monterrey}
\date{\vspace{-1.5cm}}

\begin{document}

% SLIDE 01: Portada
\begin{frame}[plain]
	\titlepage
\end{frame}

% SLIDE 02: Hoja de Ruta
\begin{frame}{Hoja de Ruta --- Capítulo 03: Variables Aleatorias Discretas}
	\vspace{-0.15cm}
	\begin{block}{Estructura Curricular de la Unidad 2}
		\footnotesize
		\begin{enumerate}
			\itemsep=0.03cm
			\item \textcolor{gray}{Sección 03.01 a 03.03: Fundamentos, PMF, CDF, Esperanza y Varianza}
			\item \textcolor{gray}{Sección 03.04: Distribuciones de Bernoulli y Binomial (Muestreo con reemplazo)}
			\item \textcolor{gray}{Sección 03.05: Distribuciones Geométrica y Binomial Negativa (Tiempos de espera)}
			\item \textbf{\textcolor{TecRojo}{Sección 03.06: Distribución Hipergeométrica}} $\leftarrow$ \emph{Foco actual}
			\item \textcolor{gray}{Sección 03.07: Distribución de Poisson y Procesos de Poisson}
		\end{enumerate}
	\end{block}
	\vspace{-0.2cm}
	\begin{alertblock}{Objetivo Didáctico}
		\scriptsize
		Dominar el modelado de extracciones aleatorias \emph{sin reemplazo} en poblaciones finitas heterogéneas, cuantificando la reducción de varianza mediante el Factor de Corrección por Población Finita (FPCF) y la covarianza negativa.
	\end{alertblock}
\end{frame}

% SLIDE 03: Motivación e Intuición
\begin{frame}{Motivación --- De la Independencia a la Dependencia Muestral}
	\vspace{-0.15cm}
	\begin{columns}[T]
		\begin{column}{0.48\textwidth}
			\begin{block}{Muestreo Con Reemplazo (Binomial)}
				\footnotesize
				\begin{itemize}
					\itemsep=0.05cm
					\item El elemento extraído retorna a la urna antes de la siguiente extracción.
					\item La probabilidad de éxito ($p$) permanece **estrictamente constante**.
					\item Los ensayos son mutuamente **independientes** ($P(A \cap B) = P(A)P(B)$).
					\item Varianza teórica máxima: $\sigma^2 = np(1-p)$.
				\end{itemize}
			\end{block}
		\end{column}
		\begin{column}{0.48\textwidth}
			\begin{alertblock}{Muestreo Sin Reemplazo (Hipergeométrica)}
				\footnotesize
				\begin{itemize}
					\itemsep=0.05cm
					\item El elemento **no retorna** a la urna; la población disminuye ($N \to N-1$).
					\item La probabilidad del siguiente intento depende del historial previo.
					\item Los ensayos muestran **dependencia correlacionada negativa**.
					\item Varianza reducida por el factor poblacional.
				\end{itemize}
			\end{alertblock}
		\end{column}
	\end{columns}
\end{frame}

% SLIDE 04: Definición Formal
\begin{frame}{Definición Formal --- La Función de Masa Hipergeométrica}
	\vspace{-0.2cm}
	\begin{block}{Estructura Poblacional ($N, K, n$)}
		\scriptsize
		En una población de tamaño $N$, exactamente $K$ elementos son \emph{éxitos} y $N-K$ son \emph{fracasos}. Al extraer una muestra de tamaño $n$ sin reemplazo, el conteo de éxitos $X$ sigue una **Distribución Hipergeométrica**: $X \sim \text{Hiper}(N, K, n)$.
	\end{block}
	\vspace{-0.25cm}
	\begin{columns}[T]
		\begin{column}{0.48\textwidth}
			\begin{exampleblock}{Fórmula Exacta de la PMF}
				\scriptsize
				Probabilidad de obtener $k$ éxitos:
				$P(X = k) = \frac{\comb{K}{k}\comb{N-K}{n-k}}{\comb{N}{n}}$.
				
				Donde $\comb{K}{k}$ elige éxitos, $\comb{N-K}{n-k}$ elige fracasos y $\comb{N}{n}$ es el total de muestras posibles.
			\end{exampleblock}
		\end{column}
		\begin{column}{0.48\textwidth}
			\begin{alertblock}{Parametrización en SciPy}
				\scriptsize
				En Python (\texttt{scipy.stats.hypergeom}):
				$\texttt{hypergeom(M, n, N)}$.
				
				Mapeo de variables:
				\texttt{M} $= N$ (Población total),
				\texttt{n} $= K$ (Éxitos totales),
				\texttt{N} $= n$ (Tamaño de muestra).
			\end{alertblock}
		\end{column}
	\end{columns}
\end{frame}

% SLIDE 05: Soporte Exacto y Vandermonde
\begin{frame}{Soporte Exacto y Normalización por Vandermonde}
	\vspace{-0.15cm}
	\begin{columns}[T]
		\begin{column}{0.48\textwidth}
			\begin{block}{Intervalo de Soporte $\mathcal{S}_X$}
				\scriptsize
				Los coeficientes exigen una doble restricción física sobre $k$:
				\begin{itemize}
					\itemsep=0.04cm
					\item Éxitos disponibles: $k \le \min(n, K)$.
					\item Fracasos disponibles: $n-k \le N-K \implies k \ge n-(N-K)$.
				\end{itemize}
				El soporte discreto exacto resulta:
				\begin{align*}
					\mathcal{S}_X = \set{\max(0, n-N+K), \dots, \min(n, K)}
				\end{align*}
			\end{block}
		\end{column}
		\begin{column}{0.48\textwidth}
			\begin{exampleblock}{Identidad de Vandermonde}
				\scriptsize
				Para garantizar que la suma total sea 1, se aplica la identidad de convolución combinatoria con $a=K$, $b=N-K$, $c=n$:
				\begin{align*}
					\sum_{k} \comb{a}{k} \comb{b}{c-k} = \comb{a+b}{c}
				\end{align*}
				\vspace{-0.25cm}
				Al dividir entre $\comb{N}{n}$, se obtiene:
				\begin{align*}
					\sum_{k \in \mathcal{S}_X} P(X=k) = 1
				\end{align*}
			\end{exampleblock}
		\end{column}
	\end{columns}
\end{frame}

% SLIDE 06: Esperanza y Linealidad
\begin{frame}{Esperanza Matemática --- Simetría y Linealidad Incondicional}
	\vspace{-0.15cm}
	\begin{block}{Descomposición en Variables Indicadoras}
		\footnotesize
		Sea $I_i = 1$ si la $i$-ésima extracción es éxito ($0$ en caso contrario), tal que $X = \sum_{i=1}^n I_i$. Por simetría marginal incondicional, cualquier elemento tiene idéntica probabilidad inicial de ser extraído en el intento $i$: $P(I_i = 1) = \frac{K}{N}$.
	\end{block}
	\vspace{-0.2cm}
	\begin{columns}[T]
		\begin{column}{0.48\textwidth}
			\begin{exampleblock}{Deducción del Primer Momento}
				\scriptsize
				Por linealidad universal de la esperanza:
				\begin{align*}
					\mathbb{E}[X] = \sum_{i=1}^n \mathbb{E}[I_i] = \sum_{i=1}^n \dfrac{K}{N} = n \dfrac{K}{N}
				\end{align*}
				¡Coincide con la media binomial ($np$) al definir $p = K/N$!
			\end{exampleblock}
		\end{column}
		\begin{column}{0.48\textwidth}
			\begin{alertblock}{Simetría y Condicionalidad}
				\scriptsize
				Aunque el ensayo 1 altera el ensayo 2 ($P(I_2=1 \mid I_1=1) = \frac{K-1}{N-1}$), el promedio total conserva la marginal:
				\begin{align*}
					P(I_2=1) = \dfrac{K}{N}\dfrac{K-1}{N-1} + \dfrac{N-K}{N}\dfrac{K}{N-1} = \dfrac{K}{N}
				\end{align*}
			\end{alertblock}
		\end{column}
	\end{columns}
\end{frame}

% SLIDE 07: Varianza y Covarianza Negativa
\begin{frame}{Covarianza Negativa y el Factor FPCF}
	\vspace{-0.15cm}
	\begin{columns}[T]
		\begin{column}{0.48\textwidth}
			\begin{block}{Covarianza Muestral ($i \neq j$)}
				\scriptsize
				La probabilidad conjunta es $P(I_i=1, I_j=1) = \frac{K(K-1)}{N(N-1)}$. Su covarianza resulta:
				\begin{align*}
					\text{Cov}(I_i, I_j) &= \dfrac{K(K-1)}{N(N-1)} - \left(\dfrac{K}{N}\right)^2 \\
					&= -\dfrac{K(N-K)}{N^2(N-1)} < 0
				\end{align*}
				Extraer un éxito reduce la probabilidad de obtener otro en el futuro.
			\end{block}
		\end{column}
		\begin{column}{0.48\textwidth}
			\begin{exampleblock}{Varianza y Corrección Poblacional}
				\scriptsize
				Sumando las varianzas individuales y las $n(n-1)$ covarianzas negativas:
				\begin{align*}
					\text{Var}(X) &= n\dfrac{K}{N}\left(1-\dfrac{K}{N}\right) \left[ \dfrac{N-n}{N-1} \right] \\
					&= np(1-p) \times \text{FPCF}
				\end{align*}
				\vspace{-0.25cm}
				El término $\text{FPCF} = \frac{N-n}{N-1} \le 1$ cuantifica la reducción teórica de incertidumbre.
			\end{exampleblock}
		\end{column}
	\end{columns}
\end{frame}

% SLIDE 08: Comparación Teórica Binomial vs Hipergeométrica
\begin{frame}{Análisis de Incertidumbre --- Censo vs. Muestreo}
	\vspace{-0.15cm}
	\begin{block}{Comportamiento Límite del FPCF ($\frac{N-n}{N-1}$)}
		\footnotesize
		El Factor de Corrección por Población Finita actúa como modulador geométrico de la varianza en función del tamaño de muestra $n$:
	\end{block}
	\vspace{-0.2cm}
	\begin{columns}[T]
		\begin{column}{0.48\textwidth}
			\begin{exampleblock}{Casos Extremos ($n=1$ y $n=N$)}
				\scriptsize
				\begin{itemize}
					\itemsep=0.04cm
					\item **Extracción Unitaria ($n=1$):** El FPCF vale $\frac{N-1}{N-1} = 1$. Sin reemplazo y con reemplazo son idénticos en un intento.
					\item **Censo Completo ($n=N$):** El FPCF vale $0 \implies \text{Var}(X) = 0$. Al medir toda la población, no existe error aleatorio ($X=K$ con certeza).
				\end{itemize}
			\end{exampleblock}
		\end{column}
		\begin{column}{0.48\textwidth}
			\begin{alertblock}{Subdispersión vs Binomial}
				\scriptsize
				Para una muestra del $40\%$ ($n = 0.40 N$ con $N=100$):
				\begin{align*}
					\text{FPCF} = \dfrac{100-40}{100-1} = \dfrac{60}{99} \approx 0.606
				\end{align*}
				La varianza muestral es un $39.4\%$ menor que si asumiéramos muestreo con reemplazo. ¡El modelo es **subdisperso** ($\sigma^2 < \mu$)!
			\end{alertblock}
		\end{column}
	\end{columns}
\end{frame}

% SLIDE 09: Convergencia Asintótica
\begin{frame}{Convergencia Asintótica --- La Regla Práctica del $5\%$}
	\vspace{-0.15cm}
	\begin{columns}[T]
		\begin{column}{0.48\textwidth}
			\begin{block}{Límite Asintótico Poblacional}
				\footnotesize
				Cuando $N \to \infty$ y $K \to \infty$ manteniendo la proporción $p = \frac{K}{N}$ fija, el FPCF tiende a 1:
				\begin{align*}
					\lim_{N \to \infty} \dfrac{N-n}{N-1} = 1
				\end{align*}
				\vspace{-0.25cm}
				Extraer sin reemplazo en una población infinita apenas altera las proporciones, y $\text{Hiper}(N, K, n) \to \text{Bin}(n, p)$.
			\end{block}
		\end{column}
		\begin{column}{0.48\textwidth}
			\begin{alertblock}{Regla de Oro en Auditorías ($\frac{n}{N} < 0.05$)}
				\footnotesize
				En ingeniería de calidad, calcular factoriales masivos en $\binom{N}{n}$ causa desbordamientos numéricos.
				
				**Criterio Práctico:** Si el tamaño muestral representa menos del **5\%** de la población ($n/N < 0.05$), se puede aproximar la Hipergeométrica mediante la Binomial con un error absoluto $< 10^{-3}$.
			\end{alertblock}
		\end{column}
	\end{columns}
\end{frame}

% SLIDE 10: Aplicaciones Clave
\begin{frame}{Aplicaciones --- Auditorías, Control y Prueba Exacta de Fisher}
	\vspace{-0.15cm}
	\begin{block}{Tres Dominios Críticos de la Distribución Hipergeométrica}
		\footnotesize
		\begin{enumerate}
			\itemsep=0.08cm
			\item **Control de Calidad (MIL-STD-105):** Inspección por atributos en lotes finitos cerrados. De un lote de $N$ piezas con $K$ defectuosas, se audita una muestra de $n$ y se rechaza si los defectos superan el límite $c$.
			\item **Auditorías Forenses y Contables:** Evaluación de padrones de facturas o cuentas transaccionales finitas para cuantificar irregularidades fiscales con rigor exacto y sin sobreestimar la varianza.
			\item **Prueba Exacta de Fisher en Bioestadística:** En ensayos clínicos $2 \times 2$ con muestras pequeñas, condicionado a totales marginales fijos, la distribución nula del número de eventos es **exactamente hipergeométrica**.
		\end{enumerate}
	\end{block}
\end{frame}

% SLIDE 11: Ejercicios Nivel Fundamental I
\begin{frame}{Ejercicios Nivel Fundamental I --- Problemas 3.6.1 y 3.6.2}
	\vspace{-0.3cm}
	\begin{columns}[T]
		\begin{column}{0.48\textwidth}
			\begin{block}{Problema 3.6.1: Soporte y Vandermonde}
				\scriptsize
				**Enunciado:** Determine el soporte de $X \sim \text{Hiper}(N, K, n)$ y verifique la suma unitaria.
				\vspace{-0.1cm}
				\begin{exampleblock}{Solución Analítica}
					Soporte: $k \le K$ y $k \ge n-(N-K) \implies$
					$\mathcal{S}_X = \set{\max(0, n-N+K), \dots, \min(n, K)}$.
					
					Identidad de Vandermonde:
					$\sum_{k} \comb{K}{k}\comb{N-K}{n-k} = \comb{N}{n} \implies \sum f_X(k) = 1$.
				\end{exampleblock}
			\end{block}
		\end{column}
		\begin{column}{0.48\textwidth}
			\begin{block}{Problema 3.6.2: Aceptación de Lotes}
				\scriptsize
				**Enunciado:** Lote $N=30$ ($K=6$ defectuosas), muestra $n=5$. Calcule $P(X \le 1)$.
				\vspace{-0.1cm}
				\begin{exampleblock}{Solución Exacta}
					Cálculo puntual ($k=0, 1$):
					$P(X=0) = \frac{\comb{6}{0}\comb{24}{5}}{\comb{30}{5}} \approx 0.2983$,
					$P(X=1) = \frac{\comb{6}{1}\comb{24}{4}}{\comb{30}{5}} \approx 0.4474$.
					
					Aceptación: $P(X \le 1) \approx 0.7457$.
				\end{exampleblock}
			\end{block}
		\end{column}
	\end{columns}
\end{frame}

% SLIDE 12: Ejercicios Nivel Fundamental II
\begin{frame}{Ejercicios Nivel Fundamental II --- Problema 3.6.3}
	\vspace{-0.15cm}
	\begin{block}{Problema 3.6.3: Linealidad Incondicional de la Esperanza}
		\footnotesize
		**Enunciado:** Sea $X = \sum_{i=1}^n I_i$ el conteo de éxitos en $n$ extracciones sin reemplazo ($I_i = 1$ si el $i$-ésimo es éxito). Demuestre por linealidad que $\mathbb{E}[X] = n \frac{K}{N}$.
	\end{block}
	\vspace{-0.2cm}
	\begin{columns}[T]
		\begin{column}{0.48\textwidth}
			\begin{exampleblock}{Simetría de la Probabilidad Marginal}
				\scriptsize
				Antes de observar los primeros ensayos, la probabilidad marginal de que el $i$-ésimo elemento sea un éxito es idéntica para toda posición $i \in \set{1, \dots, n}$:
				\begin{align*}
					P(I_i = 1) = \dfrac{K}{N} \implies \mathbb{E}[I_i] = \dfrac{K}{N}
				\end{align*}
			\end{exampleblock}
		\end{column}
		\begin{column}{0.48\textwidth}
			\begin{alertblock}{Aplicación de Linealidad Universal}
				\scriptsize
				La linealidad $\mathbb{E}[\sum Y_i] = \sum \mathbb{E}[Y_i]$ opera **sin importar la dependencia** o covarianza:
				\begin{align*}
					\mathbb{E}[X] = \sum_{i=1}^n \mathbb{E}[I_i] = \sum_{i=1}^n \dfrac{K}{N} = n \dfrac{K}{N}
				\end{align*}
				\vspace{-0.25cm}
				Demostración directa sin sumatorias algebraicas complejas.
			\end{alertblock}
		\end{column}
	\end{columns}
\end{frame}

% SLIDE 13: Ejercicios Nivel Operativo I
\begin{frame}{Ejercicios Nivel Operativo I --- Problemas 3.6.4 y 3.6.5}
	\vspace{-0.15cm}
	\begin{columns}[T]
		\begin{column}{0.48\textwidth}
			\begin{block}{Problema 3.6.4: Reducción por FPCF}
				\scriptsize
				**Enunciado:** Con $N=50, K=10, n=20$, compare la varianza Hipergeométrica vs Binomial ($p=0.20$).
				\vspace{-0.1cm}
				\begin{exampleblock}{Cálculo y Reducción Muestral}
					Varianza con reemplazo (Binomial):
					\begin{align*}
						\text{Var}_{\text{Bin}}(X) = 20(0.20)(0.80) = 3.20
					\end{align*}
					\vspace{-0.25cm}
					Factor de Corrección FPCF: $\frac{50-20}{50-1} = \frac{30}{49} \approx 0.6122$.
					\begin{align*}
						\text{Var}_{\text{Hiper}}(X) = 3.20 \times 0.6122 \approx 1.9592
					\end{align*}
					\vspace{-0.25cm}
					La varianza disminuye un $38.78\%$ por muestreo sin reemplazo.
				\end{exampleblock}
			\end{block}
		\end{column}
		\begin{column}{0.48\textwidth}
			\begin{block}{Problema 3.6.5: Auditoría Fiscal}
				\scriptsize
				**Enunciado:** En $N=100$ facturas con $K=15$ irregulares, se auditan $n=10$. Calcule $P(X \ge 3)$.
				\vspace{-0.1cm}
				\begin{exampleblock}{Cálculo por Complemento}
					Probabilidades acumuladas inferiores ($k=0,1,2$):
					\begin{align*}
						P(X=0) &= \comb{15}{0}\comb{85}{10}/\comb{100}{10} \approx 0.1868 \\
						P(X=1) &= \comb{15}{1}\comb{85}{9}/\comb{100}{10} \approx 0.3640 \\
						P(X=2) &= \comb{15}{2}\comb{85}{8}/\comb{100}{10} \approx 0.2863
					\end{align*}
					\vspace{-0.25cm}
					Suma: $P(X \le 2) \approx 0.8372 \implies P(X \ge 3) \approx 0.1628$.
				\end{exampleblock}
			\end{block}
		\end{column}
	\end{columns}
\end{frame}

% SLIDE 14: Ejercicios Nivel Operativo II
\begin{frame}{Ejercicios Nivel Operativo II --- Problema 3.6.6}
	\vspace{-0.15cm}
	\begin{block}{Problema 3.6.6: Aproximación Binomial y Error Absoluto}
		\footnotesize
		**Enunciado:** En $N=2000$ expedientes con $K=100$ errores ($p=0.05$), se auditan $n=20$. Verifique la condición $n/N < 0.05$ y calcule el error absoluto $|P(X=2) - P(Y=2)|$.
	\end{block}
	\vspace{-0.2cm}
	\begin{columns}[T]
		\begin{column}{0.48\textwidth}
			\begin{exampleblock}{Cálculo Exacto Hipergeométrico}
				\scriptsize
				Verificación: $\frac{n}{N} = \frac{20}{2000} = 0.01 = 1\% < 5\%$.
				La probabilidad exacta sin reemplazo es:
				\begin{align*}
					P(X=2) = \dfrac{\comb{100}{2}\comb{1900}{18}}{\comb{2000}{20}} \approx 0.189725
				\end{align*}
			\end{exampleblock}
		\end{column}
		\begin{column}{0.48\textwidth}
			\begin{alertblock}{Modelo Binomial y Error Absoluto}
				\scriptsize
				Aproximando con $Y \sim \text{Bin}(n=20, p=0.05)$:
				\begin{align*}
					P(Y=2) = \comb{20}{2}(0.05)^2(0.95)^{18} \approx 0.188677
				\end{align*}
				\vspace{-0.25cm}
				Error absoluto: $|0.189725 - 0.188677| = 0.001048$. ¡Una precisión superior al $99.8\%$!
			\end{alertblock}
		\end{column}
	\end{columns}
\end{frame}

% SLIDE 15: Ejercicios Nivel Analítico I
\begin{frame}{Ejercicios Nivel Analítico --- Problemas 3.6.7 y 3.6.8}
	\vspace{-0.3cm}
	\begin{columns}[T]
		\begin{column}{0.48\textwidth}
			\begin{block}{Problema 3.6.7: Momentos Factoriales}
				\scriptsize
				**Enunciado:** Deduzca $\mathbb{E}[X]$ y $\mathbb{E}[X(X-1)]$.
				\vspace{-0.15cm}
				\begin{exampleblock}{Absorción Combinatoria}
					\fontsize{7pt}{8.2pt}\selectfont
					Con $k\comb{K}{k} = K\comb{K-1}{k-1}$ y $\comb{N}{n} = \frac{N}{n}\comb{N-1}{n-1}$:
					$\mathbb{E}[X] = \sum k \frac{\comb{K}{k}\comb{N-K}{n-k}}{\comb{N}{n}} = n\frac{K}{N}$.
					
					Análogamente: $\mathbb{E}[X(X-1)] = n(n-1)\frac{K(K-1)}{N(N-1)}$.
				\end{exampleblock}
			\end{block}
		\end{column}
		\begin{column}{0.48\textwidth}
			\begin{block}{Problema 3.6.8: Varianza Exacta}
				\scriptsize
				**Enunciado:** Deduzca la varianza muestral.
				\vspace{-0.15cm}
				\begin{exampleblock}{Covarianza Muestral}
					\fontsize{7pt}{8.2pt}\selectfont
					Para $i \neq j$: $\text{Cov}(I_i, I_j) = -\frac{K(N-K)}{N^2(N-1)}$.
					
					Sumando varianzas y covarianzas:
					$\text{Var}(X) = n\frac{K}{N}(1-\frac{K}{N})(\frac{N-n}{N-1})$.
				\end{exampleblock}
			\end{block}
		\end{column}
	\end{columns}
\end{frame}

% SLIDE 16: Ejercicios Nivel Desafiante
\begin{frame}{Ejercicios Nivel Desafiante --- Problemas 3.6.9 y 3.6.10}
	\vspace{-0.15cm}
	\begin{columns}[T]
		\begin{column}{0.48\textwidth}
			\begin{block}{Problema 3.6.9: Límite Asintótico}
				\scriptsize
				**Enunciado:** Demuestre que si $N \to \infty$ con $\frac{K_N}{N} \to p$, entonces $P(X_N=k) \to \comb{n}{k}p^k(1-p)^{n-k}$.
				\vspace{-0.1cm}
				\begin{exampleblock}{Factoriales Decrecientes}
					Usando factoriales decrecientes $A^{(m)}$:
					\begin{align*}
						P(X_N=k) = \comb{n}{k} \dfrac{K_N^{(k)} (N-K_N)^{(n-k)}}{N^{(n)}}
					\end{align*}
					\vspace{-0.25cm}
					Al dividir cada factor entre $N$:
					$\frac{K_N^{(k)}}{N^k} \to p^k$, $\frac{(N-K_N)^{(n-k)}}{N^{n-k}} \to (1-p)^{n-k}$, obteniendo el límite Binomial exacto.
				\end{exampleblock}
			\end{block}
		\end{column}
		\begin{column}{0.48\textwidth}
			\begin{block}{Problema 3.6.10: Prueba Exacta de Fisher}
				\scriptsize
				**Enunciado:** En tabla $2 \times 2$ con $N=25$ ($K=12$ activos, $13$ placebo) y $n=8$ trombos, se observa $X=1$ en activo. Calcule el valor $p$.
				\vspace{-0.1cm}
				\begin{exampleblock}{Decisión Biomédica Exacta}
					Bajo $H_0$, $X \sim \text{Hiper}(25, 12, 8)$.
					Valor $p$ de cola inferior ($P(X \le 1)$):
					\begin{align*}
						P(X=0) &= \comb{12}{0}\comb{13}{8}/\comb{25}{8} \approx 0.001190 \\
						P(X=1) &= \comb{12}{1}\comb{13}{7}/\comb{25}{8} \approx 0.019039
					\end{align*}
					\vspace{-0.25cm}
					Suma: $p \approx 0.0202 < 0.05 \implies$ Rechazo de $H_0$.
				\end{exampleblock}
			\end{block}
		\end{column}
	\end{columns}
\end{frame}

% SLIDE 17: Lab Python I --- PMF y Soporte
\begin{frame}[fragile]{Laboratorio Python --- Validación Vectorizada de PMF/Soporte (`scipy.stats`)}
	\vspace{-0.2cm}
	\begin{block}{Implementación en \texttt{03.06\_hypergeometric.py} (Bloque 1)}
		\lstinputlisting[language=Python, firstline=16, lastline=33, basicstyle=\fontsize{6.2pt}{7.2pt}\ttfamily]{../../code/03_variables_aleatorias_discretas/03.06_hypergeometric.py}
	\end{block}
\end{frame}

% SLIDE 18: Lab Python II --- FPCF y Monte Carlo
\begin{frame}[fragile]{Laboratorio Python --- Covarianza Negativa y FPCF por Monte Carlo ($N=250,000$)}
	\vspace{-0.25cm}
	\begin{block}{Implementación en \texttt{03.06\_hypergeometric.py} (Bloque 2)}
		\lstinputlisting[language=Python, firstline=37, lastline=58, basicstyle=\fontsize{5.5pt}{6.5pt}\ttfamily]{../../code/03_variables_aleatorias_discretas/03.06_hypergeometric.py}
	\end{block}
\end{frame}

% SLIDE 19: Lab Python III --- Prueba de Fisher y Asintótica
\begin{frame}[fragile]{Laboratorio Python --- Prueba Exacta de Fisher y Convergencia Asintótica}
	\vspace{-0.2cm}
	\begin{block}{Implementación en \texttt{03.06\_hypergeometric.py} (Bloque 3)}
		\lstinputlisting[language=Python, firstline=62, lastline=78, basicstyle=\fontsize{6.2pt}{7.2pt}\ttfamily]{../../code/03_variables_aleatorias_discretas/03.06_hypergeometric.py}
	\end{block}
\end{frame}

% SLIDE 20: Cierre y Puente Didáctico
\begin{frame}{Cierre --- Síntesis Estocástica y Puente Didáctico}
	\vspace{-0.15cm}
	\begin{columns}[T]
		\begin{column}{0.48\textwidth}
			\begin{block}{Síntesis Conceptual de la Hipergeométrica}
				\footnotesize
				\begin{itemize}
					\itemsep=0.08cm
					\item **Muestreo Sin Reemplazo ($N, K, n$):** Modelado exacto en poblaciones finitas.
					\item **Subdispersión y FPCF:** $\text{Var}(X) = np(1-p)\left(\frac{N-n}{N-1}\right) < np(1-p)$ debido a la covarianza negativa.
					\item **Convergencia Binomial:** Si $\frac{n}{N} < 0.05$, el muestreo sin reemplazo equivale al muestreo con reemplazo.
				\end{itemize}
			\end{block}
		\end{column}
		\begin{column}{0.48\textwidth}
			\begin{alertblock}{Puente a la Sección 03.07}
				\footnotesize
				Hemos completado el estudio del conteo discreto en ensayos con $p$ fija o en poblaciones finitas con $N, K$ conocidos.
				
				¿Qué ocurre en fenómenos continuos donde los eventos suceden aleatoriamente en el **tiempo o el espacio**, con una tasa media constante de ocurrencia $\lambda$?
				
				$\to$ Esto nos conduce a la **Distribución de Poisson**.
			\end{alertblock}
		\end{column}
	\end{columns}
\end{frame}

\end{document}
